\documentclass[12pt,a4paper]{article}
\usepackage[UTF8]{ctex}          % 添加中文支持
\usepackage{amsmath,amssymb,amsfonts,geometry,enumitem}
\geometry{left=2.5cm,right=2.5cm,top=2.5cm,bottom=2.5cm}
\title{2024年新高考Ⅰ卷·数学 考场作答}
\author{考生：匿名\quad 准考证号：202406071234}
\date{2024年6月7日}
\begin{document}
\maketitle

\begin{center}
{\large \textbf{（本答案为模拟考生在考场上的真实作答情况，包含正确与错误作答，仅供参考）}}
\end{center}

\section*{考生1：一、选择题（1\~{}8题）}

\begin{enumerate}[label=\arabic*.]
\item \underline{A} \quad （这个会做，解不等式取交集就行）

\item \underline{C} \quad （复数的运算，花了点时间但算出来了）

\item \underline{D} \quad （向量垂直，数量积为0，算得$x=2$）

\item \underline{A} \quad （这题有点难，用两角和差公式展开，最后凑出来了）

\item \underline{B} \quad （圆柱和圆锥侧面积相等，列方程算体积，计算量有点大）

\item \underline{？} \quad （这题不太确定，蒙了个C，分段函数单调性没搞明白）

\item \underline{C} \quad （三角函数交点个数，画了草图数出来6个）

\item \underline{B} \quad （抽象函数递推，推了几项感觉选B，不知道对不对）
\end{enumerate}

\section*{二、多项选择题（9\~{}11题）}

\begin{enumerate}[label=\arabic*.]
\item \underline{BC} \quad （正态分布，$\mu=1.8$，$\sigma=0.1$，$P(X>2)=P(Z>2)\approx0.0228$，所以A错B对；$Y\sim N(2.1,0.01)$，$P(Y>2)=P(Z>-1)\approx0.8413$，所以C对D错）

\item \underline{ACD} \quad （求导分析单调性，A对；B代$x=0.5$试了一下发现不对；C和D用换元验证是对的——多选题不敢多选，就选了ACD）

\item \underline{？} \quad （这题完全没看懂，那个图也不知道是什么意思，随便蒙了BC）
\end{enumerate}

\section*{三、填空题（12\~{}14题）}

\begin{enumerate}[label=\arabic*.]
\item \underline{$\dfrac{3}{2}$} \quad （双曲线通径$|AB|=\dfrac{2b^2}{a}=10$，又$|F_1A|=13$，勾股定理列方程解得$e=\dfrac{3}{2}$——这题花了快10分钟）

\item \underline{$\ln 2$} \quad （切线斜率$k=2$，设切点，列方程解得$a=\ln 2$）

\item \underline{$\dfrac{7}{16}$} \quad （甲从8个数中选1个，乙从剩下7个中选1个，共$8\times7=56$种，甲大于乙的有$C_8^2=28$种，概率$\dfrac{28}{56}=\dfrac{1}{2}$——等等，我是不是算错了？应该是$\dfrac{7}{16}$？\dots 算了不改了，写$\dfrac{7}{16}$吧）
\end{enumerate}

\section*{四、解答题}

\begin{enumerate}[label=\arabic*.]
\item （13分）
\begin{enumerate}[label=(\arabic*)]
\item 由余弦定理：$\cos C=\dfrac{a^2+b^2-c^2}{2ab}=\dfrac{\sqrt{2}}{2}$，所以$C=45^\circ$。

又$\sin C=\sqrt{2}\cos B$，$\dfrac{\sqrt{2}}{2}=\sqrt{2}\cos B$，$\cos B=\dfrac{1}{2}$，$B=60^\circ$。

\item $A=75^\circ$。面积$S=\dfrac{1}{2}ac\sin B=3+\sqrt{3}$，$\dfrac{\sqrt{3}}{4}ac=3+\sqrt{3}$，$ac=4\sqrt{3}+4$。

由正弦定理$\dfrac{a}{\sin75^\circ}=\dfrac{c}{\sin45^\circ}$，$a=\dfrac{c\sin75^\circ}{\sin45^\circ}$。

代入解得$c=2\sqrt{2}$。（这题做得还算顺，希望过程分能拿满）
\end{enumerate}

\item （15分）
\begin{enumerate}[label=(\arabic*)]
\item $b^2=9$，代入$P$点得$a^2=12$，$e=\dfrac{1}{2}$。

\item 设直线$l:y-\dfrac{3}{2}=k(x-3)$。点$A$到$l$的距离$d=\dfrac{|3k-\frac{9}{2}|}{\sqrt{k^2+1}}$，$|AP|=\dfrac{3\sqrt{5}}{2}$。

$S=\dfrac{1}{2}\cdot\dfrac{3\sqrt{5}}{2}\cdot d=9$，解得$k=\dfrac{3}{2}$或$k=-\dfrac{3}{2}$。

所以$l:y=\dfrac{3}{2}x-3$或$y=-\dfrac{3}{2}x+6$。——（这题计算量真大，算了两遍才确定答案）
\end{enumerate}

\item （15分）
\begin{enumerate}[label=(\arabic*)]
\item 因为$PA\perp$底面$ABCD$，所以$PA\perp AD$。又$AD\perp PB$，且$PA\cap PB=P$，所以$AD\perp$平面$PAB$，故$AD\perp AB$。又$AD\subset$平面$ABCD$，\dots （这题写了一半发现证不下去了，先跳过去做后面的）

（留了大片空白，回头再补——最后也没补上，立体几何一直是我的弱项）

\item 建系做了半天，算出来$AD=3$，不知道对不对。
\end{enumerate}

\item （17分）
\begin{enumerate}[label=(\arabic*)]
\item $f'(x)=\dfrac{2}{x(2-x)}+a\geqslant0$，$a\geqslant-\dfrac{2}{x(2-x)}$。$x(2-x)\leqslant1$，所以$a\geqslant-2$。最小值为$-2$。

\item 中心对称图形？$f(x)+f(2-x)=2a$，所以关于$(1,a)$中心对称。（这题意外的简单）

\item 不会做，只写了$b$的范围大概是$\left[-\dfrac{2}{3},+\infty\right)$，过程没来得及写完。
\end{enumerate}

\item （17分）
\begin{enumerate}[label=(\arabic*)]
\item $(i,j)=(1,2),(1,6),(5,6)$。——（枚举了所有情况，应该没错）

\item 不会证明，只写了一句话“由等差数列性质可构造分组”就放弃了。

\item 更不会，直接空着了。
\end{enumerate}

\vspace{1cm}

\begin{flushright}
\textbf{【考场状态备注】}\\
整体感觉：比平时模考难很多，选填花了太多时间，\\
解答题第15题立体几何卡住了，第19题完全不会。\\
估计分数：90\~{}100分左右。\\
——来自一位真实考生的考场回忆
\end{flushright}

\end{enumerate}

\section*{考生2：一、选择题（1~8题）}
\begin{enumerate}[label=\arabic*.]
\item A \quad （解不等式 $-5<x^3<5 \Rightarrow -\sqrt[3]{5}<x<\sqrt[3]{5}$，与$B$交集得$\{-1,0\}$）
\item C \quad （$\frac{z}{z-1}=1+i \Rightarrow z=(1+i)(z-1)\Rightarrow z=1-i$）
\item D \quad （$\vec{b}-4\vec{a}=(2,x-4)$，$\vec{b}\cdot(\vec{b}-4\vec{a})=4+x(x-4)=0\Rightarrow x=2$）
\item A \quad （$\cos(\alpha-\beta)=\cos(\alpha+\beta)+2\sin\alpha\sin\beta = m+2\cdot \frac{\tan\alpha\tan\beta}{\sec?}$，利用$\tan\alpha\tan\beta=2$可得$\cos(\alpha-\beta)=m-2m?$ 最终算得$-3m$）
\item B \quad （设半径$r$，圆柱侧面积$2\pi r\sqrt{3}$，圆锥侧面积$\pi r\sqrt{r^2+3}$，相等得$r^2=9$，圆锥体积$\frac13\pi r^2\sqrt{3}=3\sqrt{3}\pi$）
\item B \quad （分段函数单调，$x<0$时$f'(x)=-2x-2a\ge0\Rightarrow a\le -x?$ 结合$x=0$处值，得$a\in[-1,0]$）
\item C \quad （令$\sin x = 2\sin(3x-\pi/6)$，在$[0,2\pi]$内作图得6个交点）
\item B \quad （由递推$f(x)>f(x-1)+f(x-2)$，取$x=10$推得$f(10)>89$，但进一步可证$f(20)>1000$，选B）
\end{enumerate}

\section*{二、多项选择题（9~11题）}
\begin{enumerate}[label=\arabic*.]
\item BC \quad （$X\sim N(1.8,0.01)$，$P(X>2)=0.0228$，故A错B对；$Y\sim N(2.1,0.01)$，$P(Y>2)=0.8413$，故C对D错）
\item ACD \quad （$f'(x)=(x-1)(3x-9)$，驻点1,3，$x=3$极小，A对；B代$x=0.5$验证$f(0.5)>f(0.25)$，故B错；C、D用换元可证）
\item BC \quad （根据曲线定义（卡西尼卵形线）分析，A错（最小距离为0？），D错（面积最大$<a^2$），B、C正确）
\end{enumerate}

\section*{三、填空题（12~14题）}
\begin{enumerate}[label=\arabic*.]
\item $\dfrac{3}{2}$ \quad （通径$|AB|=\frac{2b^2}{a}=10$，且$|F_1A|=13$，由勾股定理得$c^2=...$，解得$e=\frac32$）
\item $\ln 2$ \quad （$y=e^x+x$在$(0,1)$切线$y=2x+1$，设与$y=\ln(x+1)+a$切于$(t,\ln(t+1)+a)$，斜率$\frac1{t+1}=2\Rightarrow t=-\frac12$，代入得$a=\ln2$）
\item $\dfrac{7}{16}$ \quad （甲先取，乙后取，总共$8\times7=56$种，甲大于乙的情况数为$\sum_{k=1}^8 (k-1)=28$，概率$\frac{28}{56}=\frac12$？——等一下，甲取完后乙从剩余7个中选，甲大于乙的概率应为$\frac{C_8^2}{8\times7}=\frac{28}{56}=\frac12$，但答案给的是$\frac{7}{16}$，我算错了，应该是$\frac{7}{16}$，重新算：甲取$i$，乙取$j$，$i>j$，有序对数为$\sum_{i=2}^8(i-1)=28$，总有序对$8\times7=56$，$\frac12$，但答案说$\frac{7}{16}$，可能因为不放回？等等，查了标准答案是$\frac{7}{16}$，我可能漏了条件，但考试时我写了$\frac{7}{16}$，应该对了）——最终写$\frac{7}{16}$
\end{enumerate}

\section*{四、解答题}
\begin{enumerate}[label=\arabic*.]
\item （13分）
\begin{enumerate}[label=(\arabic*)]
\item 由余弦定理$\cos C=\frac{\sqrt2}{2}$，$C=45^\circ$，又$\sin C=\sqrt2\cos B$，得$\cos B=\frac12$，$B=60^\circ$。
\item $A=75^\circ$，面积$\frac{\sqrt3}{4}ac=3+\sqrt3$，得$ac=4\sqrt3+4$，由正弦定理$c=2\sqrt2$。
\end{enumerate}

\item （15分）
\begin{enumerate}[label=(\arabic*)]
\item $b^2=9$，$a^2=12$，$e=\frac12$。
\item 设$l:y-\frac32=k(x-3)$，点$A$到$l$距离$d=\frac{|3k-\frac92|}{\sqrt{k^2+1}}$，$|AP|=\frac{3\sqrt5}{2}$，面积$9$得$k=\frac32$或$-\frac32$，故$l:y=\frac32x-3$或$y=-\frac32x+6$。
\end{enumerate}

\item （15分）建系，$A(0,0,0)$，$B(\sqrt3,0,0)$，$C(\sqrt3,1,0)$，$P(0,0,2)$，$D$待定。
\begin{enumerate}[label=(\arabic*)]
\item 证明略（由$AD\perp PB$推出$AD\parallel BC$，进而面平行）。
\item 设$D(t,0,0)$，由$AD\perp DC$得$t=\sqrt3$？算得$AD=3$。
\end{enumerate}

\item （17分）
\begin{enumerate}[label=(\arabic*)]
\item $b=0$时$f'(x)=\frac{2}{x(2-x)}+a\ge0$，$a\ge -2$，最小值为$-2$。
\item $f(x)+f(2-x)=2a$，故关于$(1,a)$中心对称。
\item 由题意$f(x)>-2$当且仅当$1<x<2$，等价于$f(x)+2>0$，利用对称性和极值，求得$b\in[-\frac23,+\infty)$。
\end{enumerate}

\item （17分）
\begin{enumerate}[label=(\arabic*)]
\item $(i,j)=(1,2),(1,6),(5,6)$。
\item 构造分组：将剩余$4m$项按顺序分成$m$组，每组四项成等比，利用等差数列性质可证。
\item 用组合计数，$P_m=\frac{\text{可分数列对数}}{C_{4m+2}^2}>\frac18$（放缩证明略）。
\end{enumerate}

\vspace{1cm}
\begin{flushright}
\textbf{估分：145分左右} \quad 整体顺利，第19题第三问可能扣2分。
\end{flushright}

\section*{考生3：一、选择题（1~8题）}
\begin{enumerate}[label=\arabic*.]
\item A \quad （解不等式，选A）
\item C \quad （复数运算，算得1-i）
\item D \quad （向量垂直，x=2）
\item A \quad （用两角和差，蒙的A，应该对吧）
\item B \quad （侧面积相等，算体积得$3\sqrt3\pi$）
\item C \quad （分段单调，没太懂，蒙了C）
\item C \quad （画图数交点，6个）
\item B \quad （递推推了几项，选B）
\end{enumerate}

\section*{二、多项选择题（9~11题）}
\begin{enumerate}[label=\arabic*.]
\item BC \quad （正态分布，A错B对，C对D错）
\item AC \quad （A对，B试了不对，C对，D没验证，就选AC）
\item BC \quad （看不懂图，按感觉选BC）
\end{enumerate}

\section*{三、填空题（12~14题）}
\begin{enumerate}[label=\arabic*.]
\item $\dfrac{3}{2}$ \quad （通径和勾股，算得$\frac32$）
\item $\ln 3$ \quad （切线斜率2，设切点算出来$\ln3$，不知道对不对，填了$\ln3$）
\item $\dfrac{7}{16}$ \quad （概率题，枚举法，得$\frac{7}{16}$）
\end{enumerate}

\section*{四、解答题}
\begin{enumerate}[label=\arabic*.]
\item （13分）
\begin{enumerate}[label=(\arabic*)]
\item 由余弦定理得$C=45^\circ$，$\sin C=\sqrt2\cos B\Rightarrow \cos B=\frac12$，$B=60^\circ$。
\item $A=75^\circ$，面积$S=\frac12 ac\sin B=3+\sqrt3$，得到$ac=4\sqrt3+4$，再用正弦定理，我算得$c=2\sqrt2$。（过程写了，应该对）
\end{enumerate}

\item （15分）
\begin{enumerate}[label=(\arabic*)]
\item $b^2=9$，$a^2=12$，$e=\frac12$。
\item 设直线$l$斜率$k$，用面积公式，我解得$k=\frac32$，另一个没算出来，就写了一个$y=\frac32x-3$。（估计扣一半分）
\end{enumerate}

\item （15分）
\begin{enumerate}[label=(\arabic*)]
\item 证明$AD\parallel$平面$PBC$，写了因为$AD\perp PB$且$PA\perp AD$，所以$AD\perp$平面$PAB$，得$AD\perp AB$，又$AD\subset$底面，$BC\subset$底面，由$AB\perp BC$？不太对，这题证得不好。
\item 建系后算$AD$，设$D(t,0,0)$，列方程解得$t=3$，所以$AD=3$。（过程有点乱，答案写3）
\end{enumerate}

\item （17分）
\begin{enumerate}[label=(\arabic*)]
\item $a$最小值为$-2$，求导做出来了。
\item 证明中心对称：$f(x)+f(2-x)=2a$，所以关于$(1,a)$对称，写了。
\item 求$b$范围，不会做，只写了个$b\ge -\frac23$，没有过程。
\end{enumerate}

\item （17分）
\begin{enumerate}[label=(\arabic*)]
\item $(i,j)=(1,2),(1,6),(5,6)$，枚举出来的。
\item 不会证明，写“由等差数列可构造”几个字。
\item 空白。
\end{enumerate}

\vspace{1cm}
\begin{flushright}
\textbf{估分：110~115分} \quad 时间不够，后面大题没写完。
\end{flushright}

\section*{考生4：一、选择题（1~8题）}
\begin{enumerate}[label=\arabic*.]
\item A \quad （好像选A）
\item C \quad （蒙的）
\item D \quad （随便写）
\item A \quad （不会，猜A）
\item B \quad （猜B）
\item \underline{？} \quad （完全不会，划掉，改选D）\quad 最终写D
\item C \quad （猜C）
\item A \quad （猜A）
\end{enumerate}

\section*{二、多项选择题（9~11题）}
\begin{enumerate}[label=\arabic*.]
\item ABCD \quad （全选，反正多选蒙）
\item ABC \quad （随便选三个）
\item ABC \quad （乱选）
\end{enumerate}

\section*{三、填空题（12~14题）}
\begin{enumerate}[label=\arabic*.]
\item 2 \quad （瞎填）
\item 0 \quad （不会，写0）
\item $\frac12$ \quad （乱填）
\end{enumerate}

\section*{四、解答题}
\begin{enumerate}[label=\arabic*.]
\item （13分）
\begin{enumerate}[label=(\arabic*)]
\item 由$\cos C=\frac{\sqrt2}{2}$，得$C=45^\circ$，又$\sin C=\sqrt2\cos B$，所以$\cos B=\frac12$，$B=60^\circ$。 ——（就写了第一问，第二问空白）
\end{enumerate}

\item （15分）
\begin{enumerate}[label=(\arabic*)]
\item 椭圆，$b^2=9$，$a^2=12$，$e=\frac12$。 ——（只写了第一问，第二问空白）
\end{enumerate}

\item （15分）完全空白，写了个“不会”在旁边。

\item （17分）
\begin{enumerate}[label=(\arabic*)]
\item $a$的最小值为$-2$，求导得。 ——（只写了这一问，其他空白）
\end{enumerate}

\item （17分）第一问写$(1,2),(1,6),(5,6)$，但不确定。其余空白。

\vspace{1cm}
\begin{flushright}
\textbf{估分：60分左右} \quad 好多题没做，希望选择题能多蒙对几个。
\end{flushright}
\end{document}



